\section{Introduction}
\label{sec:intro}

Software traceability is the ability to relate each requirement to the
artifacts it derives from, satisfies, and is verified by. For airborne
software, DO-178C makes this relation a certification objective, and
certification authorities increasingly expect trace chains to be
auditable across two or more hops~\cite{do178c,arp4754a,easa2025npa,faa2024airoadmap}.
Because requirements management tools record these links explicitly, a
requirements corpus is also a typed link graph, and answering a
traceability question over it is a retrieval problem over both text and
graph.

Retrieval-augmented generation (RAG) is a natural fit for this task,
but published comparisons of RAG architectures disagree. Microsoft
GraphRAG~\cite{graphrag2024} exploits graph structure, yet loses to
vector RAG on detailed queries in recent
benchmarks~\cite{ragvsgraphrag2025,graphragbench2026}.
Adaptive-RAG~\cite{adaptiverag2024} routes queries by predicted
complexity but does not consider hop distance, and traceability systems
such as LiSSA~\cite{lissa2025} and TVR~\cite{niu_tvr2025} do not reason
over typed graphs. Each of these evaluations reports which architecture
wins in one setting, usually with one corpus, one embedder, and one
judge. They do not explain why the verdict changes when the setting
changes.

In this paper, we hold a matrix of five pipelines fixed (vanilla,
agentic, agentic with a graph tool, GraphRAG, and a learned adaptive
router) and vary four factors around it. These are the retrieval
embedder (e5-small, OpenAI text-embedding-3-small, and bge-m3, from
three vendors and at three dimensionalities), the corpus (DO-178C
requirements with typed edges, and Wikipedia paragraph chains from
MuSiQue and 2WikiMultihopQA), the graph traversal (a bounded two-hop
walk and personalized PageRank), and the faithfulness judge (GPT-5.4,
GPT-4.1, and Gemini 3.7-flash). The design comprises 13{,}456 generation
runs and more than 30{,}000 faithfulness judgments. We ask four research questions. \textbf{RQ1}: Which
pipeline produces the best answer-level citations at each hop depth,
and is that ordering stable across embedders, corpora, and seeds?
\textbf{RQ2}: Does the GraphRAG versus vector RAG comparison depend on
whether citation quality is measured on the retrieved context or on the
answer? \textbf{RQ3}: Does faithfulness decline with hop depth on every
corpus, and which property of a pipeline determines the decline?
\textbf{RQ4}: How stable are LLM faithfulness judges across time,
retrieval states, and vendors?

We find that the best pipeline depends on hop depth but not on the
embedder or the seed (RQ1); that GraphRAG ranks below the vanilla
baseline when citations are scored on the retrieved context and first
when they are scored on the answer, in all six settings we test (RQ2);
that faithfulness declines with hop depth on all three corpora once
adequately powered, and only for pipelines that expand their context
(RQ3); and that judges drift on identical inputs and reverse which of
them is stricter between corpora, while reproducing the same hop-wise
ordering (RQ4).
